\section{An example with an explicit limit}
\label{sec:example}

In this section, we give an example in which the limiting misclassification rate is
obtained in a closed form. We consider the following setting.

\begin{enumerate}[label=(E),leftmargin=*]
\item Let $n_1=n_2=1$, $\theta_1=\theta_2=\theta$ and $m_1=m_2=1$. The two
populations have a common spiked direction $h$ with
$\lambda_{11}=\lambda_{21}=\lambda=c\theta$, so that $r_{11}=1$. The spike is
orthogonal to the mean difference, that is, $h\perp\mu$ or equivalently
$\beta_{11}=\beta_{21}=0$. We have $c_{10}=c_{20}=1-c$ and $z$'s are standard
normal.
\end{enumerate}

\begin{theorem}
\label{thm:4}
Assume \textup{(E)}. Then, for $x_0\in\pi_1$, it holds that
\begin{equation}
\frac{y(x_0)}{\theta}
\wto
\frac{\delta}{2}+c(z_1-z_2)\Bigl(z_0-\frac{z_1+z_2}{2}\Bigr),
\label{eq:7.1}
\end{equation}
and, putting $A=c(z_1-z_2)$, the limiting conditional misclassification rate is
given by
\begin{equation}
e^*(1)=\Phi\Bigl(\sgn(A)\frac{z_1+z_2}{2}-\frac{\delta}{2|A|}\Bigr),
\label{eq:7.2}
\end{equation}
where $\Phi$ denotes the standard normal distribution function and $z_1,z_2$ are
i.i.d.\ standard normal.
\end{theorem}

\begin{proof}
(i) When each class has one observation, the hard-margin SVM is the perpendicular
bisector of the two points. Indeed, putting $w=x_{11}-x_{21}$ and
$b=-w^T(x_{11}+x_{21})/2$, we have $y_p(w^Tx_p+b)=1$ for $p\in\cI$, and the
maximality of the margin is easily checked. Hence
\begin{equation}
y(x_0)=(x_{11}-x_{21})^Tx_0-\frac{\|x_{11}\|^2-\|x_{21}\|^2}{2},
\label{eq:7.3}
\end{equation}
where we used $w^Tx_0+b=(x_{11}-x_{21})^Tx_0-\|x_{11}\|^2/2+\|x_{21}\|^2/2$.

(ii) We take the origin at $(\mu_1+\mu_2)/2$, so that $\mu_1=\mu/2$ and
$\mu_2=-\mu/2$. Under~(E) we have
\[
x_{11}=\tfrac{\mu}{2}+\sqrt{\lambda}\,z_1 h+\xi_1,\quad
x_{21}=-\tfrac{\mu}{2}+\sqrt{\lambda}\,z_2 h+\xi_2,\quad
x_0=\tfrac{\mu}{2}+\sqrt{\lambda}\,z_0 h+\xi_0,
\]
so that $(x_{11}-x_{21})^Tx_0=\Delta/2+\lambda(z_1-z_2)z_0+o_P(\theta)$. Using
$h\perp\mu$ and $h\perp\xi_i$, and noting that all the cross terms are $o_P(\theta)$
from Lemmas~\ref{lem:2} and~\ref{lem:3} and Step~6 of the proof of
Theorem~\ref{thm:3}, we obtain
\[
(x_{11}-x_{21})^Tx_0=\frac{\Delta}{2}+\lambda(z_1-z_2)z_0+o_P(\theta).
\]
On the other hand, from~\eqref{eq:3.2} we have
\[
\|x_{11}\|^2=\frac{\Delta}{4}+\lambda z_1^2+\tr(\Sigma_{1*})+o_P(\theta),\quad
\|x_{21}\|^2=\frac{\Delta}{4}+\lambda z_2^2+\tr(\Sigma_{2*})+o_P(\theta).
\]
Here $\tr(\Sigma_{1*})$ and $\tr(\Sigma_{2*})$ cancel out because $n_1=n_2$ and
$\theta_1=\theta_2$. We note that, in general, the bias term of the existing theory
appears at this place. Hence
\[
\frac{\|x_{11}\|^2-\|x_{21}\|^2}{2}=\frac{\lambda(z_1^2-z_2^2)}{2}+o_P(\theta).
\]
Substituting these into~\eqref{eq:7.3} and dividing by $\theta$, we obtain
\[
\frac{y(x_0)}{\theta}
=\frac{\delta}{2}+c(z_1-z_2)z_0-\frac{c(z_1^2-z_2^2)}{2}+o_P(1),
\]
and the factorization $z_1^2-z_2^2=(z_1-z_2)(z_1+z_2)$ gives~\eqref{eq:7.1}.

(iii) Conditioning on $z_1$ and $z_2$, we have $z_0\sim N(0,1)$. Let
$A=c(z_1-z_2)$. When $A>0$, we have $f^*<0$ if and only if
$z_0<\sgn(A)(z_1+z_2)/2-\delta/(2|A|)$, so that
$e^*(1)=\Phi(\sgn(A)(z_1+z_2)/2-\delta/(2|A|))$. When $A<0$, the inequality is
reversed and
\[
e^*(1)=\Phi\Bigl(\sgn(A)\frac{z_1+z_2}{2}-\frac{\delta}{2|A|}\Bigr)
\]
still holds. Combining the two cases with $\Phi(-\cdot)=1-\Phi(\cdot)$, we obtain
\eqref{eq:7.2}.
\end{proof}

\subsection{Interpretation}

The right-hand side of~\eqref{eq:7.1} is the sum of the signal $\delta/2$ coming
from the mean difference and the nearest centroid rule
$c(z_1-z_2)(z_0-(z_1+z_2)/2)$ in the one-dimensional spiked coordinate. That is,
the high dimensionality vanishes except for the spiked direction, and what remains
is a classical small-sample problem in one dimension with $N=2$ observations. We
give the following observations.

\begin{enumerate}[label=(\roman*)]
\item When $c\to 0$, which corresponds to the non-spiked case, we have $A\to 0$ and
hence $e^*(1)\to 0$ when $\delta>0$, so that $E\{e^*(1)\}\to 0$. That is, the
consistency property of the existing theory is recovered.

\item When $\delta=0$, we have $e^*(1)=\Phi(\sgn(A)(z_1+z_2)/2)$, whose expectation
is $1/2$. That is, the classification is completely random.

\item The misclassification rate increases as $c$ increases. We emphasize that the
spike degrades the performance even though it is orthogonal to the mean difference,
so that it carries no information about the classification. This phenomenon cannot
be observed in the existing theory.
\end{enumerate}
